For exercises 3.1 and 3.15, determine the range (possible values) of the random variable.

\enquestion{3.1}{
    The random variable is the number of nonconforming solder connections on a printed circuit board with $2000$ connections.
}

\enquestion{3.5}{
    A batch of $600$ machined parts contains $10$ that do not conform to customer requirements. 
    Parts are selected successively, without replacement, until a nonconforming part is obtained.
    The random variable is the number of parts selected.
}

\enquestion{3.14}{
    The sample space of a random experiment is $\{a,b,c,d,e\}$, and each outcome is equally likely.
    A random variable is defined as follows:
    \begin{center}
        \begin{tabular}{*{7}{c}}
            \toprule
                \textbf{Outcome}& $a$ & $b$ & $c$ & $d$ & $e$ & $f$ \\
            \midrule
                $x$ & $0$ & $0$ & $1.5$ & $1.5$ & $2$ & $3$ \\
            \bottomrule
        \end{tabular}
    \end{center}
    Determine the probability mass function of $X$.
    Use the probability mass function to determine the following probabilities:

    \begin{tabular}{*{2}{m{.45\textwidth}}}
        (a) $P(X=2)$ & (d) $P(0.6 < X < 2.7)$ \\
        (b) $P(X>3)$ & (e) $P(0 < X \leq 2)$ \\
        (c) $P(X=0 \text{ or } X=2)$ & \\
    \end{tabular}
}

\enquestion{3.17}{
    Verify that the following function is a probability mass function, and determine the requested probabilities.
    \[
        f(x) = \frac{2x+1}{25}, \qquad x=0,1,2,3,4
    \]
    \begin{tabular}{*{2}{m{.45\textwidth}}}
        (a) $P(X=3)$ & (b) $P(X \leq 1)$ \\
        (c) $P(2 \leq X < 4)$ & (d) $P(X > -10)$ \\
    \end{tabular}
}

\enquestion{3.29}{
    An assembly consists of three mechanical components. 
    Suppose that the probabilities that the first, second, and third components meet specifications are $0.95$, $0.98$, and $0.97$, respectively.
    Assume that the components are independent.
    Determine the probability mass function of the number of components in the assembly that meet specifications.
}

\enquestion{3.35}{
    Determine the cumulative distribution function of the random variable with the following probability mass function:
    \begin{center}
        \begin{tabular}{*{6}{c}}
            \toprule
                $x$ & $-2$ & $-1$ & $0$ & $1$ & $2$ \\
            \midrule
                $f(x)$ & $1/8$ & $2/8$ & $2/8$ & $2/8$ & $1/8$ \\
            \bottomrule        
        \end{tabular}
    \end{center}
    Furthermore, determine the following probabilities:

    \begin{tabular}{*{2}{m{.45\textwidth}}}
        (a) $P(X \leq 1.25)$ & (b) $P(X \leq 2.2)$ \\
        (c) $P(-1.1 < X \leq 1)$ & (d) $P(X > 0)$ \\
    \end{tabular}
}

\enquestion{3.43}{
    Verify that the following function is a cumulative distribution function, and determine the probability mass function and the requested probabilities:
    \[
        F(x) = \begin{cases}
            0 & x<-10\\
            0.3 & -10 \leq x < 30 \\
            0.7 & 0.7 \leq x < 50 \\
            1 & 50 \leq x
        \end{cases}
    \]
    \begin{tabular}{*{2}{m{.45\textwidth}}}
        (a) $P(X \leq 50)$ & (b) $P(X \leq 40)$ \\
        (c) $P(40 \leq X \leq 60)$ & (d) $P(X < 0)$ \\
        (e) $P(0 \leq X < 10)$ & (f) $P(-10 < X < 10)$ \\
    \end{tabular}
}


\enquestion{3.51}{
    Determine the mean and variance of the following random variable:
    \begin{center}
        \begin{tabular}{*{6}{c}}
            \toprule
                $x$ & $-2$ & $-1$ & $0$ & $1$ & $2$ \\
            \midrule
                $f(x)$ & $1/8$ & $2/8$ & $2/8$ & $2/8$ & $1/8$ \\
            \bottomrule        
        \end{tabular}
    \end{center}
}


\enquestion{3.53}{
    Determine the mean and variance of the following random variable with the following probability mass function:
    \[
        f(x) = \frac{2x+1}{25}, \qquad x=0,1,2,3,4
    \]
}


\enquestion{3.67}{
    Let the random variable $X$ have a discrete uniform distribution on the integers $1 \leq x \leq 5$.
    Determine the mean and variance of $X$.
}


\enquestion{3.68}{
    Thickness measurements of a coating process are made to the nearest hundredth of a millimeter.
    The thickness measurements are uniformly distributed with values $0.14$, $0.15$, $0.16$, $0.17$, $0.18$, $0.19$, and $0.20$.
    Determine the mean and variance of the coating thickness for this process.
}

\enquestion{3.70}{
    The lengths of plate glass parts are measured to the nearest tenth of a millimeter.
    The lengths are uniformly distributed with values at every tenth of a millimeter starting at $640.0$ and continuing through $640.9$.
    Determine the mean and variance of the lengths.
}